\section{Methods}
\label{sec:method}

% Theory, methods, algorithms. Assume the reader has familiarity with basic
% linear algebra and statistics. Break into \subsection{}s as the content
% develops, e.g. one per model/algorithm.

This section presents the models, ... We assume the reader has familiarity
with basic linear algebra and statistics.

\subsection{Logistic regression}\label{subsec:logistic-regression}

Logistic regression is a natural choice as a baseline. It is a simple model in which
everything can be computed, its cost function is convex, and it is exactly a neural
network with no hidden layers~\cite{hjorth-jensen_2026}.

\subsubsection{Error measures}

We evaluate a prediction $\bm{\tilde{y}}$ of targets $\bm{y}$ by

\subsection{Neural networks}\label{subsec:neural-networks}



\subsubsection{The artificial neuron}\label{subsubsec:artificial-neuron}


\subsubsection{Activation functions}

% Introduce why it must be non-linear and why we care about the derivative
% Element-wise (done for all $z$ in $\bm{z}$). Should we write z or \bm{z} below?

Our study will revolve around the following three activation functions:

\paragraph{Sigmoid} The logistic or sigmoid function is given by
\begin{equation*}\label{eq:sigmoid}
    \sigma(z) = \frac{1}{1+e^{-z}},
\end{equation*}


\paragraph{ReLU} The rectified linear unit (ReLU),
\begin{equation*}\label{eq:relu}
    \mathrm{ReLU}(z) = \max (0, z),
\end{equation*}
is faster to compute

The leaky ReLU is defined as
\begin{equation*}\label{eq:leaky-relu}
    \mathrm{LReLU}(z) =
    \begin{cases}
        z, & \text{if } z > 0, \\
        \alpha z, & \text{if } z \leq 0,
    \end{cases}
    \qquad \alpha \approx 0.01.
\end{equation*}
